% !TeX root = ../main.tex
\section{Beispieldatensatz für das Fundsachensystem}
Die folgende Tabelle enthält zufällig zusammengestellte Beispieldaten zur
Veranschaulichung der Erfassung und Verwaltung von Fundgegenständen. Sie stellt
keine tatsächlichen Meldungen oder Evaluationsergebnisse dar.

\begin{table}[htbp]
    \centering
    \caption{Beispielhafte Fundgegenstände und Bearbeitungsstatus}
    \label{tab:fundgegenstaende-beispiel}
    \begin{tabular}{llll}
        \toprule
        Gegenstand & Kategorie & Fundort & Status \\
        \midrule
        Schwarze Trinkflasche & Trinkflasche & Turnsaal & Abholbereit \\
        Blaues Federpennal & Schulsachen & Raum B204 & Nicht zugeordnet \\
        Grauer Regenschirm & Kleidung/Zubehör & Haupteingang & Abgeholt \\
        Roter USB-Stick & Elektronik & Bibliothek & In Verwahrung \\
        Grüne Sporttasche & Tasche & Garderobe & Abholbereit \\
        \bottomrule
    \end{tabular}
\end{table}
\chapter{Ergänzende Materialien}
Hier können Testprotokolle, Fragebögen oder ergänzende Angaben stehen.
